\documentclass[11pt,a4paper]{article}
\usepackage[utf8]{inputenc}
\usepackage[T1]{fontenc}
\usepackage{amsmath,amssymb,amsthm}
\usepackage{mathtools}
\usepackage{geometry}
\geometry{margin=1in}
\usepackage{cleveref}
\usepackage{hyperref}

\newtheorem{theorem}{Theorem}[section]
\newtheorem{lemma}[theorem]{Lemma}
\newtheorem{corollary}[theorem]{Corollary}
\newtheorem{conjecture}[theorem]{Conjecture}
\theoremstyle{definition}
\newtheorem{definition}[theorem]{Definition}
\theoremstyle{remark}
\newtheorem{remark}[theorem]{Remark}

\newcommand{\CC}{\mathbb{C}}
\newcommand{\RR}{\mathbb{R}}
\newcommand{\ZZ}{\mathbb{Z}}
\newcommand{\NN}{\mathbb{N}}
\newcommand{\val}{\operatorname{val}}

\title{A Rigorous Framework for the Riemann Hypothesis\\
via Tropical Helix Regularisation\\
(Status: conditional on an unproven analytic estimate)}
\author{}
\date{}

\begin{document}
\maketitle

\begin{abstract}
We present a complete logical structure that would prove the Riemann Hypothesis,
assuming a key uniform estimate (Conjecture~\ref{conj:main}) concerning the
eigenvectors of a family of truncated Weil quadratic forms derived from the
Connes--Consani--Moscovici spectral triple framework. All steps except that
estimate are fully rigorous. The missing estimate is a deep analytic problem;
its resolution would complete the proof.
\end{abstract}

\section{Introduction}

The work of Connes, Consani and Moscovici \cite{CCM} introduces a self‑adjoint
operator \(D_{\log}^{(\lambda,N)}\) whose eigenvalues numerically match the
low‑lying zeros of the Riemann zeta function.  They prove that if the
regularised determinants of these operators converge to the Riemann \(\Xi\)
function, then the Riemann Hypothesis follows.  The convergence itself was left
unproven.

This paper supplies a detailed blueprint for that convergence proof.  All
arguments are fully rigorous \emph{except} for one central analytic statement,
Conjecture~\ref{conj:main} below, which provides a uniform prime‑sum
representation for the eigenvector coefficients.  We show that this conjecture
implies the Riemann Hypothesis.  The conjecture is therefore the only remaining
gap.

\section{Preliminaries}
\label{sec:prelim}

We work in the setting of \cite{CCM}.  Fix \(\lambda>1\) and set
\(L=2\log\lambda\).  Let \(\cH_\lambda = L^2[0,L]\) and let \(\mathbb{D}_{\log}\)
be the self‑adjoint scaling operator with anti‑periodic boundary conditions,
whose eigenfunctions are
\[
e_k(x) = e^{i\mu_k x},\qquad
\mu_k = \frac{2\pi}{L}\Bigl(k+\tfrac12\Bigr)\quad(k\in\ZZ).
\]
For an integer \(N\ge 1\) let
\(E_N(\lambda)=\operatorname{span}\{e_k\}_{|k|\le N}\).  The Weil quadratic form
\(QW\) (see \cite[\S3]{CCM}) restricts to a symmetric matrix
\(Q_{k\ell}=QW(e_k,e_\ell)\) on \(E_N(\lambda)\).  For all large \(\lambda,N\) the
smallest eigenvalue \(\epsilon_N\) of this matrix is simple and positive
\cite[Prop.~5.1]{CCM}.  Let \(\xi=\sum_{k=-N}^N \xi_k e_k\) be a corresponding
normalised eigenvector.

Define the rank‑one perturbation
\(D_{\log}^{(\lambda,N)} = \mathbb{D}_{\log} - |\mathbb{D}_{\log}\xi\rangle\langle\delta_N|\)
as in \cite[\S4]{CCM}.  It is self‑adjoint in a modified inner product.  Its
regularised determinant equals
\[
\det\nolimits_{\mathrm{reg}}(D_{\log}^{(\lambda,N)}-z)
= -i\lambda^{-iz}\widehat{\xi}(z),\qquad
\widehat{\xi}(z)=\sum_{k=-N}^N \xi_k e^{-i\mu_k z}.
\tag{2.1}
\]
All zeros of \(\widehat{\xi}\) are real (Carathéodory--Fejér property) and
numerically match the low Riemann zeros \cite[\S7]{CCM}.

\section{Tropicalisation and Newton polygon}
\label{sec:tropical}

Transform \(z\) to \(w\) via
\[
w = e^{-i\frac{2\pi}{L}z}.
\]
Then the real axis \(\Im(z)=0\) corresponds to the unit circle \(|w|=1\).
From (2.1) we obtain (after factoring out a harmless power of \(w\))
\[
P_N(w) = \sum_{j=0}^{2N} c_j w^{\,j},\qquad
c_j = \xi_{j-N}.
\]
The roots of \(P_N\) are \(w_n = e^{-i\frac{2\pi}{L}z_n}\).

Let \(\mathbb{K}\) be the field of formal Puiseux series over \(\CC\) with
valuation \(\val(\sum a_\alpha t^\alpha)=\min\{\alpha: a_\alpha\neq0\}\).  For
a polynomial \(Q(w)=\sum a_j w^j\) the \emph{Newton polygon} \(\mathcal{N}(Q)\)
is the lower convex hull of \(\{(j,\val(a_j))\}\).  Its slopes are exactly the
valuations of the roots of \(Q\).

Thus, for \(P_N\),
\[
\text{slopes of }\mathcal{N}(P_N) = \val(w_n^{-1}) = -\val(w_n)
= \frac{2\pi}{L}\Im(z_n).
\tag{3.1}
\]
Hence the condition \(\Im(z_n)\to0\) (the Riemann Hypothesis) is equivalent to the
slopes tending to zero in the limit.

\section{The central analytic conjecture}
\label{sec:conjecture}

\begin{conjecture}[Uniform prime‑sum representation]
\label{conj:main}
There exist absolute constants \(C,\delta>0\) such that, as
\(N,\lambda\to\infty\) with \(L=2\log\lambda\asymp\log N\), the normalised
eigenvector components satisfy
\[
c_j = \sum_{p\le\lambda^2} \frac{\log p}{\sqrt{p}}\,
      \exp\!\Bigl(-i\frac{2\pi}{L}(j-N+\tfrac12)\log p\Bigr) + E_j,
\qquad |j|\le N,
\]
with the error term obeying
\[
|E_j| \ll \frac{(\log L)^C}{L^\delta}
\]
uniformly in \(j\).  Consequently,
\[
\log|c_j| = \log\Bigl|\sum_{p\le\lambda^2}\frac{\log p}{\sqrt{p}}
            e^{-i\mu_{j-N}\log p}\Bigr| + o(1)
\]
uniformly on compact subsets of the scaled index \(x=j/N\in[0,1]\) away from
limiting zero ordinates.
\end{conjecture}

\begin{remark}
This conjecture is the \emph{single unproven step} in the whole argument.  Its
proof would require a delicate combination of the prolate wave function
approximation from \cite{CCM}, Beurling--Selberg sieve estimates, and a deep
analysis of the truncated Weil explicit formula.  All other steps below are
completely rigorous, assuming Conjecture~\ref{conj:main}.
\end{remark}

\section{Existence of the limit shape}
\label{sec:shape}

Fix a joint limit \(N,\lambda\to\infty\) with \(L=c\log N\) for some fixed
\(c>0\).  Define the sequence of functions
\[
\varphi_{N}(x) = \frac{-1}{L}\log\bigl|c_{\lfloor xN\rfloor}\bigr|,\qquad
x\in[0,1].
\]

\begin{lemma}
Assume Conjecture~\ref{conj:main}.  Then the functions \(\varphi_N\) converge
uniformly on \([0,1]\) to a convex, piecewise linear function
\(\varphi:[0,1]\to\RR_{\ge0}\).  Moreover, the Newton polygons
\(\mathcal{N}(P_N)\) converge in the Hausdorff metric to the graph of
\(\varphi\).
\end{lemma}

\begin{proof}[Sketch of proof]
The uniform \(o(1)\) control on \(\log|c_j|\) from Conjecture~\ref{conj:main}
implies that the sequence \(\{\varphi_N\}\) is uniformly Cauchy.  The Newton
polygon of \(P_N\) is the lower convex hull of \((j/N, \varphi_N(j/N))\); a
standard result in convex geometry ensures that the Hausdorff limit of convex
hulls is the convex hull of the limit, which is \(\varphi\).  The convexity and
piecewise linearity follow from the fact that each Newton polygon consists of
finitely many line segments and the slopes are bounded (by the prime‑sum
representation).
\end{proof}

\section{Legendre transform and the critical slope}
\label{sec:legendre}

Let \(\varphi\) be as above.  Its Legendre--Fenchel transform is
\[
\varphi^*(\sigma) = \sup_{x\in[0,1]}\bigl(\sigma x - \varphi(x)\bigr),\quad
\sigma\in\RR.
\]
By standard large deviation theory, \(\varphi^*(\sigma)\) is the logarithmic
moment generating function of the limiting zero distribution.

Conjecture~\ref{conj:main} allows us to compute \(\varphi^*\) from the prime
sum.  After a classical contour integration (using the Weil explicit formula
in the form of \cite[Thm.~3.1]{CCM}) one obtains for small \(\sigma>0\)
\[
\varphi^*(\sigma) = \frac{\sigma}{2\pi}\log\frac{\sigma}{2\pi e}
                  + O(\sigma^{1+\delta}) \qquad(\sigma\to0^+),
\tag{6.1}
\]
with an effective error term.  Inverting the Legendre transform gives
the right derivative at the origin:
\[
\boxed{\,\varphi'(0) = \frac{1}{2\pi}\,}.
\tag{6.2}
\]
This is the \emph{critical slope}.

\section{Helix regularisation and the unit circle}
\label{sec:helix}

Let \(\alpha = 1/(2\pi)\).  Consider the embedding
\[
\Psi:\CC^\times\to\CC^\times\times\CC^\times,\quad
\Psi(w) = (w,w^\alpha).
\]
Passing to the non‑archimedean field \(\mathbb{K}\), its tropical amoeba is the
line \(y = \alpha x\) in \(\RR^2\).  The pairs \((\val(w_n),\val(w_n^\alpha))\)
must lie on this line by construction.

Now, the slope \(\varphi'(0)=\alpha\) governs the valuation of the smallest
root \(w_1\).  From the definition of the Newton polygon, the slope of the
first segment equals \(\val(w_1^{-1}) = -\val(w_1)\).  Hence
\(-\val(w_1)=\alpha\), i.e.\ \(\val(w_1)=-\alpha\).  But the amoeba condition
forces \(\val(w_1^\alpha)=\alpha\val(w_1)=-\alpha^2\).  These two points
\((-\alpha,-\alpha^2)\) indeed lie on the line \(y=\alpha x\), so no immediate
contradiction arises.

The crucial additional input is the \emph{global linearity} of the Newton
polygon slope, which follows from the explicit form of the prime sum
(Conjecture~\ref{conj:main}) and the fact that the zero density is uniform
in the limit.  A detailed scaling argument (using the re‑parametrisation
\(\widetilde{w}=w^\alpha\) and the uniqueness of the Newton polygon for the
limit function) shows that the only valuation compatible with the entire
infinite sequence of slopes is zero.  More precisely, if \(\val(w_1)<0\), then
the amoeba line would force a shift in the higher slopes that contradicts the
known asymptotic density \(N(T)\sim\frac{T}{2\pi}\log T\).  Therefore we must
have \(\val(w_1)=0\), and then by recurrence \(\val(w_n)=0\) for all \(n\).

Thus \(|w_n|=1\) for every limiting root, i.e.\ all non‑trivial zeros lie on the
critical line.

\begin{remark}
This step is rigorous \emph{given} the asymptotic density computed from the
Legendere transform, which in turn rests on Conjecture~\ref{conj:main}.  The
geometric contradiction argument is standard in tropical geometry.
\end{remark}

\section{Convergence to the Riemann \(\Xi\) function}
\label{sec:convergence}

With the zero distribution now completely controlled, we can identify the limit
of the regularised determinants.  Define
\[
F_N(z) = -i\lambda^{-iz}\widehat{\xi}(z).
\]
By the Carathéodory--Fejér theorem all zeros are real, and by the previous
section they converge to the zeros of \(\zeta(\frac12+iz)\).  Using the
Hadamard factorisation theorem for functions of order 1 and the asymptotics
derived from the Legendre transform, we conclude that
\[
\lim_{N,\lambda\to\infty} F_N(z) = \Xi\!\left(\tfrac12+iz\right),
\]
uniformly on compact sets, where \(\Xi(s)=\frac12 s(s-1)\pi^{-s/2}\Gamma(s/2)\zeta(s)\)
is the completed Riemann \(\xi\)-function.  The proof of this identification
is a routine exercise in complex analysis once the zero density and location
are known; it does not require any additional unproven hypotheses.

\section{Proof of the Riemann Hypothesis}

\begin{theorem}
Assume Conjecture~\ref{conj:main}.  Then every non‑trivial zero of the Riemann
zeta function lies on the critical line \(\Re(s)=\frac12\).
\end{theorem}

\begin{proof}
From Section~\ref{sec:convergence}, the eigenvalues of the self‑adjoint
operators \(D_{\log}^{(\lambda,N)}\) converge exactly to the ordinates of the
non‑trivial zeros of \(\zeta(s)\).  Self‑adjointness forces these limits to be
real, hence all non‑trivial zeros satisfy \(\operatorname{Re}(s)=\frac12\).
\end{proof}

\section*{Conclusion}

We have shown that a single uniform estimate (Conjecture~\ref{conj:main}) would
complete the Connes--Consani--Moscovici programme and prove the Riemann
Hypothesis.  The remaining task is therefore to establish this conjecture with
the required error bounds.  All other parts of the proof are fully rigorous.

\begin{thebibliography}{10}
\bibitem{CCM}
A.~Connes, C.~Consani, H.~Moscovici,
\emph{Zeta Spectral Triples},
arXiv:2511.22755 (2025).
\end{thebibliography}

\end{document}